\documentclass[11pt,a4paper]{article}

\usepackage[T2A]{fontenc}
\usepackage[utf8]{inputenc}
\usepackage[russian]{babel}
\usepackage[a4paper,hmargin=1.5cm,vmargin=1.5cm]{geometry}
\usepackage{amsmath,amssymb}
\usepackage{array,booktabs}
\usepackage{eso-pic}

\setlength{\parindent}{0pt}
\setlength{\parskip}{4pt}
\pagestyle{empty}

\newcommand{\ych}{y_{\text{чн}}}

% Карточка варианта: номер, уравнение для вида y_чн, уравнение для общего решения
% Размеры карточки: сетка 2x2 на всю страницу с зазором \cardgap для разрезания
\newlength{\cardgap}\setlength{\cardgap}{0.4cm}
\newlength{\cardw}\newlength{\cardh}
\newcommand{\card}[3]{%
  {\fboxrule=0pt\fbox{\begin{minipage}[t][\cardh][t]{\cardw}
    \vspace{4pt}
    \textbf{Вариант #1}

    \smallskip
    Фамилия: \rule{5cm}{0.4pt}

    \medskip
    \textbf{1.} \ Записать вид частного решения уравнения (коэффициенты не находить):
    \[
      #2
    \]

    \medskip
    \textbf{2.} \ Найти общее решение уравнения
    \[
      #3
    \]
  \end{minipage}}}}

\begin{document}

% Страница с карточками: поля 5 мм (меньше принтер обычно не пропечатывает)
\newgeometry{margin=0.5cm}
\AddToShipoutPictureBG*{%
  \AtPageCenter{\makebox(0,0){\rule{0.4pt}{\dimexpr\paperheight-1cm\relax}}}%
  \AtPageCenter{\makebox(0,0){\rule{\dimexpr\paperwidth-1cm\relax}{0.4pt}}}}
\setlength{\cardw}{\dimexpr(\textwidth-\cardgap)/2-2\fboxsep-2\fboxrule\relax}
\setlength{\cardh}{\dimexpr(\textheight-\cardgap)/2-2\fboxsep-2\fboxrule-2pt\relax}

% ---------------------------------------------------------------------
% Страница 1: варианты для печати и разрезания
{\setlength{\parskip}{0pt}\offinterlineskip
\noindent\card{1}{y'' - 4y = e^{2x} + x}{y'' - 3y' + 2y = 6e^{-x}}\hfill
\card{2}{y'' + 9y = \cos 3x + e^{x}}{y'' - y' - 2y = 3e^{2x}}\par

\vspace{\cardgap}

\noindent\card{3}{y'' - 2y' + y = x e^{x} + \sin x}{y'' + 2y' + y = 2\cos x}\hfill
\card{4}{y'' + y' = x^2 + e^{-x}}{y'' - 2y' = 4x}\par}

\newpage
\restoregeometry
% ---------------------------------------------------------------------
% Страница 2: ответы для преподавателя
\begin{center}
{\Large\bfseries Квиз: ответы}
\end{center}

Проверяет домашнее задание к \texttt{lesson\_03.tex}: вид частного решения по правой части и метод неопределённых коэффициентов.

В задаче 1 правая часть "--- сумма двух слагаемых, для одного из них $\lambda$ "--- корень характеристического уравнения: нужна суперпозиция и множитель $x^k$. В задаче 2 система для коэффициентов не больше $2 \times 2$. Типы задач по вариантам:
\begin{center}
\renewcommand{\arraystretch}{1.3}
\begin{tabular}{@{}lll@{}}
\toprule
Вариант & Задача 1 & Задача 2\\
\midrule
1 & $\lambda = 2$, $k = 1$ \ + \ многочлен, $k = 0$ & $\lambda = -1$, $k = 0$ \\
2 & $\lambda = 3i$, $k = 1$ \ + \ экспонента, $k = 0$ & $\lambda = 2$, $k = 1$ \\
3 & $\lambda = 1$, $k = 2$ \ + \ тригонометрия, $k = 0$ & тригонометрия, $\lambda = i$, $k = 0$ \\
4 & многочлен, $k = 1$ \ + \ $\lambda = -1$, $k = 1$ & корень $\lambda_1 = 0$, многочлен, $k = 1$ \\
\bottomrule
\end{tabular}
\end{center}

\par\medskip\textbf{Вариант 1}\par

\textbf{1.} $\lambda_{1,2} = \pm 2$. \ $e^{2x}$: $\lambda = 2$, $k = 1$; \ $x$: $\lambda = 0$, $k = 0$ \quad$\Rightarrow$\quad $\ych = Ax e^{2x} + Bx + C$

\textbf{2.} $\lambda_1 = 1$, $\lambda_2 = 2$. \ $\lambda = -1$, $k = 0$, \ $\ych = Ae^{-x}$: \ $(1 + 3 + 2)A = 6$, \ $A = 1$ \quad$\Rightarrow$\quad $y = C_1 e^{x} + C_2 e^{2x} + e^{-x}$

\par\medskip\textbf{Вариант 2}\par

\textbf{1.} $\lambda_{1,2} = \pm 3i$. \ $\cos 3x$: $\lambda = 3i$, $k = 1$; \ $e^{x}$: $\lambda = 1$, $k = 0$ \quad$\Rightarrow$\quad $\ych = x(A\cos 3x + B\sin 3x) + Ce^{x}$

\textbf{2.} $\lambda_1 = -1$, $\lambda_2 = 2$. \ $\lambda = 2$, $k = 1$, \ $\ych = Axe^{2x}$, \ $u = Ax$: \ $(u'' + 3u')\,e^{2x} = 3Ae^{2x}$, \ $A = 1$

\quad$\Rightarrow$\quad $y = C_1 e^{-x} + C_2 e^{2x} + x e^{2x}$

\par\medskip\textbf{Вариант 3}\par

\textbf{1.} $\lambda_1 = \lambda_2 = 1$. \ $xe^{x}$: $\lambda = 1$, $k = 2$; \ $\sin x$: $\lambda = i$, $k = 0$ \quad$\Rightarrow$\quad $\ych = x^2(Ax + B)\, e^{x} + C\cos x + D\sin x$

\textbf{2.} $\lambda_1 = \lambda_2 = -1$. \ $\lambda = i$, $k = 0$, \ $\ych = A\cos x + B\sin x$:
\ $\ych'' + 2\ych' + \ych = 2B\cos x - 2A\sin x = 2\cos x$, \ $A = 0$, $B = 1$

\quad$\Rightarrow$\quad $y = (C_1 + C_2 x)\, e^{-x} + \sin x$

\par\medskip\textbf{Вариант 4}\par

\textbf{1.} $\lambda_1 = 0$, $\lambda_2 = -1$. \ $x^2$: $\lambda = 0$, $k = 1$; \ $e^{-x}$: $\lambda = -1$, $k = 1$ \quad$\Rightarrow$\quad $\ych = x(Ax^2 + Bx + C) + Dxe^{-x}$

\textbf{2.} $\lambda_1 = 0$, $\lambda_2 = 2$. \ $\lambda = 0$, $k = 1$, \ $\ych = x(Ax + B)$: \ $2A - 2(2Ax + B) = 4x$, \ $A = -1$, $B = -1$

\quad$\Rightarrow$\quad $y = C_1 + C_2 e^{2x} - x^2 - x$

\end{document}
